\section{AIME I 1990}
        \begin{enumerate}
        	\item (\textit{AIME I 1990}) The \href{https://artofproblemsolving.com/wiki/index.php?title=Increasing\_sequence}{increasing sequence} $2,3,5,6,7,10,11,\ldots$ consists of all \href{https://artofproblemsolving.com/wiki/index.php?title=Positive\_integer}{positive integers} that are neither the \href{https://artofproblemsolving.com/wiki/index.php?title=Perfect\_square}{ square} nor the \href{https://artofproblemsolving.com/wiki/index.php?title=Perfect\_cube}{ cube} of a positive integer. Find the 500th term of this sequence.


 

	\item (\textit{AIME I 1990}) Find the value of $(52+6\sqrt{43})^{3/2}-(52-6\sqrt{43})^{3/2}$.


 

	\item (\textit{AIME I 1990}) Let $P_1^{}$ be a regular $r~\mbox{gon}$ and $P_2^{}$ be a regular $s~\mbox{gon}$ $(r\geq s\geq 3)$ such that each interior angle of $P_1^{}$ is $\frac{59}{58}$ as large as each interior angle of $P_2^{}$. What's the largest possible value of $s_{}^{}$?


 

	\item (\textit{AIME I 1990}) Find the positive solution to


 $\frac 1{x^2-10x-29}+\frac1{x^2-10x-45}-\frac 2{x^2-10x-69}=0$

	\item (\textit{AIME I 1990}) Let $n^{}_{}$ be the smallest positive integer that is a multiple of $75_{}^{}$ and has exactly $75_{}^{}$ positive integral divisors, including $1_{}^{}$ and itself. Find $\frac{n}{75}$.


 

	\item (\textit{AIME I 1990}) A biologist wants to calculate the number of fish in a lake. On May 1 she catches a random sample of 60 fish, tags them, and releases them. On September 1 she catches a random sample of 70 fish and finds that 3 of them are tagged. To calculate the number of fish in the lake on May 1, she assumes that 25\% of these fish are no longer in the lake on September 1 (because of death and emigrations), that 40\% of the fish were not in the lake May 1 (because of births and immigrations), and that the number of untagged fish and tagged fish in the September 1 sample are representative of the total population. What does the biologist calculate for the number of fish in the lake on May 1? 


 

	\item (\textit{AIME I 1990}) A triangle has vertices $P_{}^{}=(-8,5)$, $Q_{}^{}=(-15,-19)$, and $R_{}^{}=(1,-7)$. The equation of the bisector of $\angle P$ can be written in the form $ax+2y+c=0_{}^{}$. Find $a+c_{}^{}$.


 

	\item (\textit{AIME I 1990}) In a shooting match, eight clay targets are arranged in two hanging columns of three targets each and one column of two targets. A marksman is to break all the targets according to the following rules: 


 1) The marksman first chooses a column from which a target is to be broken. 


 2) The marksman must then break the lowest remaining target in the chosen column. 


 If the rules are followed, in how many different orders can the eight targets be broken?


 
                    \begin{figure}[H]
                        \centering
                        \includegraphics[width=0.5\textwidth]{Problem/AIME_1990_I_8_Problem_Figure_1.png}
                    \end{figure}
                    

	\item (\textit{AIME I 1990}) A fair coin is to be tossed $10_{}^{}$ times. Let $i/j^{}_{}$, in lowest terms, be the probability that heads never occur on consecutive tosses. Find $i+j_{}^{}$. 


 

	\item (\textit{AIME I 1990}) The sets $A = \{z : z^{18} = 1\}$ and $B = \{w : w^{48} = 1\}$ are both sets of complex roots of unity.  The set $C = \{zw : z \in A ~ \mbox{and} ~ w \in B\}$ is also a set of complex roots of unity.  How many distinct elements are in $C^{}_{}$?


 

	\item (\textit{AIME I 1990}) Someone observed that $6! = 8 \cdot 9 \cdot 10$.  Find the largest positive integer $n^{}_{}$ for which $n^{}_{}!$ can be expressed as the product of $n - 3_{}^{}$ consecutive positive integers.


 

	\item (\textit{AIME I 1990}) A regular 12-gon is inscribed in a circle of radius 12.  The sum of the lengths of all sides and diagonals of the 12-gon can be written in the form


 $a + b \sqrt{2} + c \sqrt{3} + d \sqrt{6},$where $a^{}_{}$, $b^{}_{}$, $c^{}_{}$, and $d^{}_{}$ are positive integers.  Find $a + b + c + d^{}_{}$.


 

	\item (\textit{AIME I 1990}) Let $T = \{9^k : k ~ \mbox{is an integer}, 0 \le k \le 4000\}$.  Given that $9^{4000}_{}$ has 3817 digits and that its first (leftmost) digit is 9, how many elements of $T_{}^{}$ have 9 as their leftmost digit?


 

	\item (\textit{AIME I 1990}) The rectangle $ABCD^{}_{}$ below has dimensions $AB^{}_{} = 12 \sqrt{3}$ and $BC^{}_{} = 13 \sqrt{3}$.  Diagonals $\overline{AC}$ and $\overline{BD}$ intersect at $P^{}_{}$.  If triangle $ABP^{}_{}$ is cut out and removed, edges $\overline{AP}$ and $\overline{BP}$ are joined, and the figure is then creased along segments $\overline{CP}$ and $\overline{DP}$, we obtain a triangular pyramid, all four of whose faces are isosceles triangles.  Find the volume of this pyramid.


 
                    \begin{figure}[H]
                        \centering
                        \includegraphics[width=0.5\textwidth]{Problem/AIME_1990_I_14_Problem_Figure_1.png}
                    \end{figure}
                    


 

	\item (\textit{AIME I 1990}) Find $ax^5 + by^5$ if the real numbers $a,b,x,$ and $y$ satisfy the equations

 \begin{align*} ax + by &= 3, \\ ax^2 + by^2 &= 7, \\ ax^3 + by^3 &= 16, \\ ax^4 + by^4 &= 42. \end{align*}



 

\end{enumerate}